%%
%% NIE-PML semestral project paper template (ACM sigconf).
%%
\documentclass[sigconf]{acmart}

\usepackage{amsmath}
\usepackage{amsfonts}
\usepackage{array}

%% \BibTeX command to typeset BibTeX logo in the docs
\AtBeginDocument{%
  \providecommand\BibTeX{{%
    Bib\TeX}}}

%% Remove ACM rights and reference information (not needed for the course).
\setcopyright{none}
\settopmatter{printacmref=false}
\renewcommand\footnotetextcopyrightpermission[1]{}
\acmConference[NIE-PML]{NIE-PML Semestral Project}{2026}{Prague, Czech Republic}

\begin{document}

\title{Positive-Unlabeled Structure Learning for Graph Bottlenecked Recommendation}

\author{Robert Bakhtigariev}
\email{bakhtrob@fel.cvut.cz}
\affiliation{%
  \institution{Faculty of Electrical Engineering, Czech Technical University in Prague}
  \city{Prague}
  \country{Czech Republic}
}

\begin{abstract}
Graph-based social recommenders rely on user-user ties that are both noisy and incomplete. Graph Bottlenecked Social Recommendation (GBSR) can only prune observed ties, and the rewiring method SHaRe adds and removes ties with fixed similarity rules. We extend GBSR to bidirectional refinement with PU-GIB, which treats the observed social graph as the output of a two-sided noisy channel, with label frequency $c$ and false-positive rate $\varepsilon$, applied to a latent preference-consistent graph. The channel is estimated from interaction-only user similarities with a noise-corrected Elkan-Noto estimator tightened by a cross-fitted anchor, and we prove that the resulting prior can only err in a permissive direction whenever the assumed noise rate does not exceed the true one. PU-GIB keeps GBSR's backbone and objective, extends its edge mask to candidate ties propagated in a separate channel, and anchors the mask to this prior with a Kullback-Leibler term. We evaluate PU-GIB against GBSR and SHaRe under equal training and a tuned weight of the social channel, on two real datasets and on a synthetic benchmark with a known latent graph, which measures directly how well each method identifies false ties and recovers missing ones.
\end{abstract}

\keywords{Social Recommendation, Graph Structure Learning, Variational Bayesian Inference, Positive-Unlabeled Learning, Graph Neural Networks}

\maketitle

\section{Introduction}
Social ties are a common auxiliary signal for mitigating interaction sparsity in collaborative filtering: grounded in homophily~\cite{mcpherson2001birds}, social Graph Neural Networks (GNNs) aggregate representations over multi-hop social neighborhoods~\cite{10.1145/3308558.3313488, 9311623}. Observed social graphs, however, are imperfect realizations of preference alignment. They are \emph{noisy} ($\varepsilon > 0$): users connect for reasons unrelated to consumption, and aggregating over such ties degrades accuracy~\cite{10.1145/3543507.3583374, 10.1145/3637528.3671807, 10.1145/3696410.3714874}. They are also \emph{incomplete} ($c < 1$): many preference-sharing pairs are simply not linked~\cite{10.1609/aaai.v39i17.33966}.

Graph Bottlenecked Social Recommendation (GBSR)~\cite{10.1145/3637528.3671807} addresses noise with a learned mask over observed ties, trained with the recommendation loss and a Hilbert-Schmidt Independence Criterion (HSIC)~\cite{gretton2005measuring} penalty in the spirit of the Information Bottleneck~\cite{tishby2000informationbottleneckmethod}. Since the mask covers observed ties only, GBSR implicitly assumes completeness ($c = 1$). SHaRe~\cite{10.1145/3589334.3645460} adds and removes links, but with fixed similarity rules (Section~\ref{sec:share}). Extending GBSR to add ties raises two problems. First, Positive-Unlabeled (PU) link recovery~\cite{10.1145/1401890.1401920, 10.1609/aaai.v39i17.33966}, the natural source of new ties, assumes that observed ties are clean ($\varepsilon = 0$), which contradicts the denoising premise. Second, GBSR's objective gives its mask no information about which unobserved pairs are plausible ties: HSIC rewards refined representations that differ from those of the observed graph, and with candidate ties in the support this pressure alone does not distinguish plausible from implausible additions.

This work extends GBSR. We propose \textbf{PU-GIB}, which keeps GBSR's backbone, mask, and objective, and adds what is missing for bidirectional refinement: a pool of plausible candidate ties, ranked by interaction similarity, and a prior that states how strongly a candidate should be believed relative to an observed tie. Within each group, the mask distinguishes individual ties through its input features and the recommendation loss. We model the observed graph as a latent preference-consistent graph $Z$ passed through a two-sided channel $(c, \varepsilon)$, estimate the channel from interaction-only similarities, and turn it by Bayes' rule into an asymmetric prior $p(Z \mid A)$ that enters the objective as a Kullback-Leibler (KL) term on the mask. Candidates are ranked by interaction similarity, as in SHaRe; the novelty lies not in this ordering but in how strongly observed and candidate ties are believed a priori, which is estimated rather than fixed by a rule, and in combining this belief with recommendation evidence in one posterior.

Evaluating structure refinement fairly requires care, because a refined graph can change accuracy without selecting better ties: a mask that rescales all ties changes the overall weight of the social channel, and a representation regularizer can help regardless of which ties are kept. We therefore train all methods identically, tune a global social weight for every social method, and include controls for both effects. Our contributions are:
\begin{itemize}
    \item \textbf{A two-sided channel prior for social graphs} that governs pruning and addition jointly and contains pruning-only denoising ($c \to 1^-$) and clean-positive PU completion ($\varepsilon \to 0^+$) as limits.
    \item \textbf{Leakage-free channel estimation with a known error direction:} a noise-corrected Elkan-Noto estimator on interaction-only similarities, tightened by a cross-fitted anchor, with a proof that the estimated prior is permissive, never restrictive, for any assumed noise rate not exceeding the true one.
    \item \textbf{A prior-anchored extension of GBSR} that applies GBSR's mask to observed ties and to candidate ties, propagated in a separate channel, and adds the prior as a KL term to GBSR's objective.
    \item \textbf{An evaluation protocol that separates tie selection from reweighting and regularization:} comparison with GBSR and SHaRe under equal training and tuned social weight, controls with uniform reweighting and with an untrained GBSR mask, a decomposition into pruning and addition, ground-truth measurements on synthetic data, and an analysis of which users benefit.
\end{itemize}

\section{Related Work}

\subsection{Social Recommendation and Denoising}
GNN-based collaborative filtering~\cite{10.1145/3331184.3331267, 10.1145/3397271.3401063} improves on neighborhood, factorization, and linear autoencoder models~\cite{10.1145/371920.372071, 5197422, 10.1145/3308558.3313710, 10.1145/3523227.3551482} by capturing high-order connectivity. Social GNNs model influence diffusion over the user-user graph~\cite{10.1145/3308558.3313488, 9311623}, and self-supervised variants add hypergraph or triadic views~\cite{yu2021self, yu2021socially}. Since observed ties often contradict preferences, denoising methods sample consistent neighbors~\cite{10.1145/3404835.3463028}, generate reliable friends adversarially~\cite{8970896, 9239885}, prune by preference-guided confidence (GDMSR)~\cite{10.1145/3543507.3583374}, or use self-augmentation~\cite{wang2023denoised}, hierarchical denoising~\cite{10771708}, and diffusion~\cite{10.1145/3696410.3714874}; earlier work used social regularization and autoencoders~\cite{10.1145/1935826.1935877, Pan2017ACD}. Graph IB methods learn structures with non-informative priors: GIB~\cite{wu2020graph} samples neighborhoods from the observed topology, and VIB-GSL~\cite{sun2022vib} learns a graph over node pairs with a prior that does not depend on whether a pair is observed. GBSR~\cite{10.1145/3637528.3671807}, our base, propagates over a joint user-item and user-user adjacency, gates observed social links with a learned mask, and minimizes HSIC between representations propagated on the masked and on the observed graph; it assigns zero mass outside the observed ties and cannot recover missing ones.

\subsection{Social Graph Rewiring and SHaRe}
\label{sec:share}
SHaRe~\cite{10.1145/3589334.3645460} is the closest method to ours, since it also removes and adds social links. It computes the cosine similarity $c_{ij}$ of user embeddings from an encoder run on the interaction graph only, removes every social link with $c_{ij} \le 0$, adds the $M$ most similar unlinked pairs, where $M$ equals the number of removed links, and weights the remaining links by normalized similarity; an auxiliary contrastive loss calibrates the embeddings. PU-GIB uses the same kind of interaction-only similarity but differs in three respects (Table~\ref{tab:comparison}). (i) \emph{Who decides a tie:} in SHaRe a fixed similarity rule; in PU-GIB a posterior probability learned from the recommendation loss. (ii) \emph{How much observation counts:} SHaRe's asymmetry between observed and missing ties is fixed by its rules for every graph; PU-GIB derives it from the estimated channel through the prior $\rho^{\text{obs}} > \rho^{\text{un}}$. (iii) \emph{How much to add:} SHaRe adds as many links as it removes; PU-GIB learns the retention of observed ties and of candidates separately.

\begin{table}[t]
\caption{How PU-GIB differs from its base GBSR and from SHaRe.}
\label{tab:comparison}
\small
\setlength{\tabcolsep}{3pt}
\begin{tabular}{@{}>{\raggedright\arraybackslash}p{0.25\columnwidth}>{\raggedright\arraybackslash}p{0.2\columnwidth}>{\raggedright\arraybackslash}p{0.22\columnwidth}>{\raggedright\arraybackslash}p{0.25\columnwidth}@{}}
\toprule
 & GBSR & SHaRe & PU-GIB \\
\midrule
Removes ties & learned mask & rule $c_{ij} \le 0$ & learned mask \\
Adds ties & no & top-$M$, $M$ = \#cuts & two-hop pool, learned \\
Mask trained by & BPR + HSIC & --- & BPR + HSIC + prior KL \\
Observed vs.\ missing & observed only & fixed by rules & estimated prior \\
\bottomrule
\end{tabular}
\end{table}

\subsection{Positive-Unlabeled Learning}
PU learning treats unlabeled instances as a mixture of hidden positives and negatives~\cite{10.1145/1401890.1401920, bekker2020learning, kiryo2017positive}. The Elkan-Noto estimator of the label frequency assumes class separability and underestimates it under overlap~\cite{bekker2020learning, lazecka2021estimating}. PU ideas have been applied to user-item feedback~\cite{hsieh2015pu, saito2020unbiased}, to node classification~\cite{10.1145/3450316, 9679099}, and to edge-incomplete link prediction in PULL~\cite{10.1609/aaai.v39i17.33966}. These methods assume clean positives and labeling that is Selected Completely at Random (SCAR) rather than dependent on topology (SAR~\cite{bekker2019beyond}); we keep PU estimation on preference features only and move topological evidence into the posterior.

\section{Method}
PU-GIB proceeds in two stages. Before training, it defines a population of candidate pairs, estimates the observation channel on it, and derives the prior (Sections~\ref{sec:universe}--\ref{sec:estimation}). During training, it learns GBSR's mask over observed and candidate ties with the prior as an additional KL term (Section~\ref{sec:vbi}).

\subsection{Two-Hop Universe and Channel Model}
\label{sec:universe}
Let $R_{\text{train}} \in \{0,1\}^{M \times N}$ be the training interactions and $A$ the symmetrized social adjacency with edge set $\mathcal{E}_{\text{obs}}$. Before any graph construction we hold out a random $10\%$ of the social edges, $\mathcal{E}_{\text{val}}$, leaving training edges $\mathcal{E}_{\text{train}}^{\text{S}}$ with adjacency $A^{\text{train}}$ and neighborhoods $\mathcal{N}(i)$; all compared methods use the same $\mathcal{E}_{\text{train}}^{\text{S}}$. We work in the universe of pairs, linked or not, that share a common neighbor:
\begin{equation}
    \Omega = \big\{\{i,j\} \;\big|\; i \neq j, \; \mathcal{N}(i) \cap \mathcal{N}(j) \neq \emptyset\big\}. \label{eq:omega_def}
\end{equation}
Every pair in $\Omega$ is labeled by $A^{\text{train}}$; held-out edges $\mathcal{P}_{\text{val}} = \mathcal{E}_{\text{val}} \cap \Omega$ stay in $\Omega$ with label $0$, as any model trained on $A^{\text{train}}$ sees them. The training labels are then a uniform thinning of the observed ones, and all channel parameters below refer to $A^{\text{train}}$: the thinning multiplies $c$, $\varepsilon$, and $p_A$ by $0.9$ and leaves $\pi$ and $\rho^{\text{obs}}$ unchanged. We split $\mathcal{P}_{\text{val}}$ at random into two halves, $\mathcal{P}_{\text{est}}$, used only by the estimator, and $\mathcal{P}_{\text{rec}}$, used only to measure the recovery of missing ties. Candidates are drawn from $\mathcal{U}_{\text{rank}} = \Omega \setminus \mathcal{E}_{\text{train}}^{\text{S}}$, which contains $\mathcal{P}_{\text{val}}$.

For each pair we posit a latent $Z_{ij} \in \{0,1\}$, with $Z_{ij}=1$ for a preference-consistent tie and prior $\pi = P(Z_{ij}=1)$, observed through a two-sided channel that we treat as SCAR within $\Omega$:
\begin{equation}
    P(A_{ij}^{\text{train}}=1 \mid Z_{ij}=1) = c, \qquad P(A_{ij}^{\text{train}}=1 \mid Z_{ij}=0) = \varepsilon, \label{eq:channel}
\end{equation}
with $0 < \varepsilon < c < 1$. $Z$ is never observed: the estimator below identifies it through preference similarity, and the posterior refines it through the recommendation loss. The observed rate $p_A = |\mathcal{E}_{\text{train}}^{\text{S}} \cap \Omega| / |\Omega| = \pi c + (1-\pi)\varepsilon$ is known. Bayes' rule gives a prior $\rho_{ij} = p(Z_{ij}=1 \mid A_{ij}^{\text{train}})$ with two values,
\begin{equation}
    \rho^{\text{obs}} = \frac{\pi c}{p_A}, \qquad \rho^{\text{un}} = \frac{\pi(1-c)}{1 - p_A}, \label{eq:rho}
\end{equation}
for observed and unobserved pairs, with $\rho^{\text{un}} < \pi < \rho^{\text{obs}}$. As $c \to 1^-$, $\rho^{\text{un}} \to 0^+$ and the support collapses to observed links, as in pruning-only denoising~\cite{10.1145/3637528.3671807}; as $\varepsilon \to 0^+$, $\rho^{\text{obs}} \to 1^-$ and all observed links are retained, as in clean-positive PU completion~\cite{10.1609/aaai.v39i17.33966}. Training edges outside $\Omega$ (bridges) keep the prior $\rho^{\text{obs}}$, an extrapolation we state explicitly.

\subsection{Leakage-Free Channel Estimation}
\label{sec:estimation}
Embeddings trained with social message passing would align linked users regardless of preferences, and topological indices predict links independently of preferences; either would let an observation classifier exploit artifacts rather than $Z$. We therefore take the embeddings $\tilde{\mathbf{e}}$ of a LightGCN~\cite{10.1145/3397271.3401063} trained on interactions only (best validation checkpoint) and use the single feature $s_{ij} = \cos(\tilde{\mathbf{e}}_i, \tilde{\mathbf{e}}_j)$, assuming
\begin{equation}
    A_{ij}^{\text{train}} \perp s_{ij} \mid Z_{ij}. \label{eq:scar_feature}
\end{equation}
This assumption can fail if very similar users find each other more easily, or if existing ties cause shared consumption (social contagion). Both push the estimate below upward, against the direction of Proposition~\ref{prop:bracket}, so the guarantee holds only under~\eqref{eq:scar_feature}; the resulting net bias cannot be measured without the latent graph, which is why we also evaluate the estimator on synthetic data where $Z$ is known.

\textbf{Elkan-Noto estimate.} Following Elkan and Noto~\cite{10.1145/1401890.1401920}, we fit a non-decreasing scorer $g(s) \approx P(A^{\text{train}}=1 \mid s)$ on $\Omega$ by isotonic regression with the Pool-Adjacent-Violators algorithm~\cite{ayer1955empirical, barlow1972statistical, zadrozny2002transforming}, which minimizes the binary cross-entropy among monotone functions~\cite{brummer2013pav}, and average it on held-out positives: $\bar{g}_{\text{val}} = |\mathcal{P}_{\text{est}}|^{-1} \sum_{\mathcal{P}_{\text{est}}} g(s_{ij})$. Held-out pairs enter the fit with label $0$ but form a small fraction of $\Omega$, so their effect on $g$ is negligible. Under class separability, $\mathbb{E}[\bar{g}_{\text{val}}] = \rho^{\text{obs}} c + (1-\rho^{\text{obs}})\varepsilon$; inverting with~\eqref{eq:rho} gives, for an assumed noise rate $\varepsilon \in (0, p_A)$,
\begin{equation}
    \hat{c}_{\text{EN}}(\varepsilon) = \frac{p_A(\bar{g}_{\text{val}} - \varepsilon)}{p_A - \varepsilon}, \label{eq:noise_corrected_en}
\end{equation}
which reduces to the classical $\hat{c} = \bar{g}_{\text{val}}$ as $\varepsilon \to 0$.

\textbf{Cross-fitted anchor.} Under~\eqref{eq:scar_feature}, $g(s) = c\,p(s) + \varepsilon(1-p(s)) \le c$ for every $s$, with $p(s) = P(Z=1 \mid s)$, so the largest value of $g$ is also a lower bound on $c$, and it is not affected by class overlap in the bulk of the data. The top block of an isotonic fit is biased upward on finite samples, so we cross-fit: we split $\Omega$ into two random halves, select the top block (merged down to at least 1000 pairs) on one half, measure the observation rate within its similarity range on the other half, and average both directions to obtain $\hat{g}_{\text{top}}$. The final estimates are
\begin{equation}
    \hat{c}(\varepsilon) = \max\{\hat{c}_{\text{EN}}(\varepsilon), \hat{g}_{\text{top}}\}, \qquad \hat{\pi}(\varepsilon) = \frac{p_A - \varepsilon}{\hat{c}(\varepsilon) - \varepsilon}, \label{eq:c_hat}
\end{equation}
and plugging them into~\eqref{eq:rho} yields the estimated prior $\hat{\rho}$. In practice the classes overlap and $\varepsilon$ is a hyperparameter; the following result fixes the direction of the resulting error.

\begin{proposition}
\label{prop:bracket}
Assume~\eqref{eq:channel} with true parameters $0 < \pi < 1$ and $0 < \varepsilon^\ast < c$, assumption~\eqref{eq:scar_feature}, a calibrated $g$, and uniformly held-out positives, and let $\bar{g}_{\text{val}}$ denote its expectation. Then, for every assumed $0 < \varepsilon \le \varepsilon^\ast$ with $\hat{c}(\varepsilon) > p_A$: (i) $\hat{c}(\varepsilon) \le c$; (ii) $\hat{\pi} \ge \pi$, $\hat{\rho}^{\text{obs}} \ge \rho^{\text{obs}}$, and $\hat{\rho}^{\text{un}} \ge \rho^{\text{un}}$.
\end{proposition}

\begin{proof}[Proof sketch]
By~\eqref{eq:scar_feature}, $g(s) = c\,p(s) + \varepsilon^\ast(1-p(s)) \le c$, so $\hat{g}_{\text{top}} \le c$, and $\mathbb{E}[\bar{g}_{\text{val}}] = cm + \varepsilon^\ast(1-m)$ with $m = \mathbb{E}[p(s) \mid A^{\text{train}}=1]$. Since $P(Z=1 \mid s, A^{\text{train}}=1) = c\,p(s)/g(s) \ge p(s)$, averaging gives $\rho^{\text{obs}} \ge m$, with equality iff $p(s) \in \{0,1\}$; as~\eqref{eq:noise_corrected_en} is increasing in $\bar{g}_{\text{val}}$, $\hat{c}_{\text{EN}}(\varepsilon^\ast) \le c$. Further, $\hat{c}_{\text{EN}}(\varepsilon) - p_A = p_A(\bar{g}_{\text{val}} - p_A)/(p_A - \varepsilon)$. If $\bar{g}_{\text{val}} \le p_A$, then $\hat{c}_{\text{EN}}(\varepsilon) \le p_A < c$ directly. Otherwise $\partial \hat{c}_{\text{EN}}/\partial \varepsilon = p_A(\bar{g}_{\text{val}} - p_A)/(p_A-\varepsilon)^2 > 0$, so $\hat{c}_{\text{EN}}(\varepsilon) \le \hat{c}_{\text{EN}}(\varepsilon^\ast) \le c$, which together with $\hat{g}_{\text{top}} \le c$ proves (i). For (ii), $\hat{\pi}$, $\hat{\rho}^{\text{obs}} = \hat{c}\hat{\pi}/p_A$, and $\hat{\rho}^{\text{un}} = \hat{\pi}(1-\hat{c})/(1-p_A)$ are decreasing in both $\hat{c}$ and $\varepsilon$ on $\varepsilon < p_A < \hat{c}$, and equal the true values at $(\varepsilon^\ast, c)$.
\end{proof}

Two consequences guide the design. Underestimating the noise rate preserves the direction of the error, whereas overestimating it can break the bound, so $\kappa = \varepsilon/p_A$ is tuned on a deliberately low grid. And the estimated prior is \emph{permissive}: it overstates the prior belief in both observed ties and candidates, so the KL term can never over-prune relative to the true channel, and pruning of noisy ties must come from the recommendation loss. The anchor can only reduce this error, because both estimators are lower bounds on $c$. At the population level $\bar{g}_{\text{val}} = \mathbb{E}[g^2]/p_A \ge p_A$ by Jensen's inequality, with equality only for a constant $g$, so the fallback below can only be triggered by an uninformative feature or by sampling noise.

\textbf{Admissibility and fallback.} The prior is valid, $0 < \hat{\rho}^{\text{un}} < \hat{\rho}^{\text{obs}} < 1$, iff $p_A < \hat{c} < 1$. Values of $\varepsilon$ for which $\hat{c}_{\text{EN}}(\varepsilon) \ge 1$, that is, $\bar{g}_{\text{val}} \ge 1 - \varepsilon(1-p_A)/p_A$, are removed from the grid, and $\hat{c}$ is clipped below 1 for numerical safety. If $\bar{g}_{\text{val}}$ does not exceed $p_A$ significantly (one-sided $z$-test at level $0.05$, using the standard error of the mean of $g$ over $\mathcal{P}_{\text{est}}$), the feature is uninformative and the prior falls back to an \emph{uninformed channel prior} computed from $p_A$ alone, by substituting $c^{\text{u}} = (1+p_A)/2$ and $\varepsilon^{\text{u}} = \kappa p_A$ into~\eqref{eq:c_hat} for $\pi$ and then into~\eqref{eq:rho}; ranking is unchanged, so a fallback run coincides with the PU-GIB-Fix ablation.

\textbf{Candidate pool.} Under~\eqref{eq:scar_feature}, the posterior probability of a hidden tie among unobserved pairs increases with $s_{ij}$, so we rank by $s_{ij}$ and form the symmetric union $\mathcal{E}_{\text{cand}}$ of each user's top-$K$ pairs in $\mathcal{U}_{\text{rank}}$. We write $\mathcal{E}_{\text{cand}}(u) = \{v \mid \{u,v\} \in \mathcal{E}_{\text{cand}}\}$ for the candidates of user $u$; users without training ties have no two-hop pairs and receive no candidates. The frozen support of the mask is $\mathcal{C} = \mathcal{E}_{\text{train}}^{\text{S}} \cup \mathcal{E}_{\text{cand}}$, with $|\mathcal{E}_{\text{cand}}| \le MK$.

\subsection{Prior-Anchored Mask Learning}
\label{sec:vbi}
\textbf{Mask.} As in GBSR, a mask assigns each tie in $\mathcal{C}$ a retention probability $\alpha_{ij}$. It reads the frozen interaction-only embeddings $\tilde{\mathbf{e}}$, so that retained ties cannot make their endpoints look more similar to the mask through propagation, together with the observation indicator and the Resource Allocation index~\cite{zhou2009predicting}, $r_{ij} = \sum_{k} |\mathcal{N}(k)|^{-1}$ over common neighbors $k \in \mathcal{N}(i) \cap \mathcal{N}(j)$; topology is admissible here because the posterior, unlike the estimator, does not rely on SCAR:
\begin{equation}
\begin{split}
    \alpha_{ij} = \sigma\big(\text{MLP}_\phi[\, & \tilde{\mathbf{e}}_i + \tilde{\mathbf{e}}_j,\; |\tilde{\mathbf{e}}_i - \tilde{\mathbf{e}}_j|,\; \tilde{\mathbf{e}}_i \odot \tilde{\mathbf{e}}_j, \\
    & \log(1+r_{ij}),\; A_{ij}^{\text{train}}\,]\big),
\end{split} \label{eq:sym_mlp}
\end{equation}
where the brackets denote concatenation. All inputs are symmetric, so $\alpha_{ij} = \alpha_{ji}$. In training we sample Binary Concrete weights~\cite{maddison2017concrete, jang2017categorical} $\tilde{Z}_{ij} = \tilde{Z}_{ji} = \sigma\big((\operatorname{logit}\alpha_{ij} + \operatorname{logit} u_{ij})/\tau\big)$ with $u_{ij} = u_{ji} \sim \text{U}(0,1)$ and temperature $\tau$; at test time each tie is weighted by $\alpha_{ij}$.

\textbf{Propagation.} Following GBSR, user-item edges and observed ties form one adjacency normalized with joint degrees $d_u = |\mathcal{N}_{\text{I}}(u)| + |\mathcal{N}(u)|$ and $d_i = |\mathcal{N}_{\text{I}}(i)|$, where $\mathcal{N}_{\text{I}}(\cdot)$ denotes interaction neighbors. Candidates form a separate channel normalized with their own degrees $d^{\text{c}}_u = |\mathcal{E}_{\text{cand}}(u)|$, so that adding them does not change the normalization of a user's interactions and observed ties:
\begin{align}
    \mathbf{e}_u^{(l+1)} = {} & \sum_{i \in \mathcal{N}_{\text{I}}(u)} \frac{\mathbf{e}_i^{(l)}}{\sqrt{d_u d_i}} + \omega \sum_{v \in \mathcal{N}(u)} \frac{\tilde{Z}_{uv}\, \mathbf{e}_v^{(l)}}{\sqrt{d_u d_v}} \notag \\
    & + \omega_{\text{c}} \sum_{v \in \mathcal{E}_{\text{cand}}(u)} \frac{\tilde{Z}_{uv}\, \mathbf{e}_v^{(l)}}{\sqrt{d^{\text{c}}_u d^{\text{c}}_v}}, \label{eq:lgcn_u} \\
    \mathbf{e}_i^{(l+1)} = {} & \sum_{u \in \mathcal{N}_{\text{I}}(i)} \frac{\mathbf{e}_u^{(l)}}{\sqrt{d_i d_u}}. \label{eq:lgcn_i}
\end{align}
The candidate term is zero for users without candidates. Final embeddings average layers $0, \dots, L$ as in LightGCN. All degrees are \emph{static}: the mask multiplies already normalized entries, so uniformly down-weighting all ties of a user is not cancelled by renormalization. The scalars $\omega$ and $\omega_{\text{c}}$ weight the two social channels; $\omega = 1$ without candidates recovers GBSR's backbone.

\textbf{Objective.} On $\mathcal{C}$ we approximate the empirical-Bayes posterior $p(Z \mid A^{\text{train}}, R_{\text{train}})$, which is proportional to $p(R_{\text{train}} \mid Z)\, p(Z \mid A^{\text{train}})$, by the factorized distribution $q_\phi(Z) = \prod_{\mathcal{C}} \text{Bern}(Z_{ij}; \alpha_{ij})$, using the BPR loss~\cite{rendle2009bpr} as a surrogate for $-\ln p(R_{\text{train}} \mid Z)$. Adding this KL term to GBSR's objective gives
\begin{equation}
\begin{split}
    \mathcal{L} = {} & \frac{1}{|\mathcal{B}|}\mathbb{E}_{\tilde{Z}}\big[\mathcal{L}_{\text{BPR}}(\mathcal{B}; \tilde{Z}, \theta)\big] + \lambda\, \mathcal{H}_\sigma(\mathcal{B}) \\
    & + \frac{\beta}{|\mathcal{C}|} \sum_{\{i,j\} \in \mathcal{C}} \mathrm{KL}\big(\text{Bern}(\alpha_{ij}) \,\|\, \text{Bern}(\hat{\rho}_{ij})\big),
\end{split} \label{eq:objective}
\end{equation}
where, as in GBSR, $\mathcal{H}_\sigma(\mathcal{B})$ is the sum of the HSIC between the $\ell_2$-normalized representations of the batch users propagated on the masked graph and on the observed graph (the first two terms of~\eqref{eq:lgcn_u} with $\tilde{Z} \equiv 1$) and the same quantity for the batch items, with a Gaussian kernel of width $\sigma$; the expectation is over relaxed samples $\tilde{Z}$; the KL divergence between Bernoulli distributions is evaluated in closed form. With $\beta = 0$ and no candidates, \eqref{eq:objective} is GBSR's objective. The KL term keeps each tie at the belief implied by the channel unless the recommendation loss provides evidence against it. Pairing a relaxed likelihood with a discrete KL makes~\eqref{eq:objective} a surrogate that recovers the discrete $\beta$-ELBO as $\tau \to 0$.

\textbf{Training.} All model parameters are trained from scratch, with the same optimizer and budget as the baselines; the interaction-only LightGCN provides only the mask's input features and the similarity $s_{ij}$. The temperature $\tau$ is annealed over the full training budget, so the mask is trained throughout. Building $\Omega$ costs $O(\sum_k |\mathcal{N}(k)|^2)$ once; propagation costs $O((|R_{\text{train}}| + |\mathcal{C}|)d)$ per epoch, plus the HSIC term on each batch.

\section{Experiments}

\subsection{Research Questions}
\textbf{RQ1}: Does PU-GIB improve top-$N$ recommendation over GBSR, SHaRe, and the controls under equal training and tuned social weight? \textbf{RQ2}: Does PU-GIB recover the latent structure, that is, how accurate is the estimated channel, how well does the mask separate false from genuine ties, and how many missing ties enter the candidate pool and the posterior? \textbf{RQ3}: Where do differences come from: pruning versus addition, calibrated versus uninformed prior, learned tie selection versus HSIC regularization and uniform reweighting, and which users benefit? \textbf{RQ4}: How robust are the methods to injected false ties?

\subsection{Datasets}
We use two public social recommendation datasets and one synthetic dataset (Table~\ref{tab:data}). \textbf{LastFM} (HetRec 2011) contains each user's listened artists, capped at 50 per user, and mutual friendships; listening counts are binarized. \textbf{Douban-Book} is taken with the preprocessing of the GBSR authors~\cite{10.1145/3637528.3671807} and split with our protocol; directed relations are symmetrized. Our numbers are therefore not directly comparable with those reported in the GBSR paper.

The \textbf{synthetic} dataset makes the latent graph observable. Users and items receive random taste vectors of dimension 16, and each user interacts with about 30 items sampled with probability increasing in user-item taste similarity. The latent graph $Z$ connects each user to the 10 users with the most similar tastes. The observed graph keeps each edge of $Z$ with probability $c = 0.5$ and adds uniformly random false edges amounting to $20\%$ of the observed edges; for RQ4 we also use $50\%$ and $70\%$. All generator parameters are fixed before any method is run.

\begin{table}[t]
\caption{Dataset statistics (before splitting). Social edges are undirected. Synthetic: range over the three generated worlds (seeds 0--2).}
\label{tab:data}
\small
\setlength{\tabcolsep}{4pt}
\begin{tabular}{@{}lrrrr@{}}
\toprule
Dataset & Users & Items & Interactions & Social edges \\
\midrule
LastFM & 1,892 & 17,632 & 92,834 & 12,717 \\
Douban-Book & 12,859 & 22,294 & 598,420 & 83,246 \\
Synthetic & 2,000 & 2,000 & 59,637--59,939 & 7,285--7,445 \\
\bottomrule
\end{tabular}
\end{table}

\subsection{Protocol and Metrics}
\label{sec:protocol}
\textbf{Splits and accuracy.} Interactions are split 80/10/10 per user into training, validation, and test sets, and $10\%$ of social edges are held out as $\mathcal{E}_{\text{val}}$; all results are averaged over three seeds of both splits. We report Recall@$N$ and NDCG@$N$ for $N \in \{10, 20\}$ under full ranking over all items not seen in training, averaged over users with held-out items. Model selection and early stopping use validation NDCG@20 for every method, including GBSR. Hyperparameters are selected on seed 0 and then fixed for all seeds. Because seeds share splits across methods, we report paired differences with respect to LightGCN-S.

\textbf{Fairness controls.} All methods share the backbone~\eqref{eq:lgcn_u}--\eqref{eq:lgcn_i}, the optimizer, and the training budget. The social weights $\omega$ and $\omega_{\text{c}}$ are tuned once per dataset on LightGCN-S and LightGCN-PU and then shared by all social methods, so no method gains from a differently weighted social channel; methods that add ties, including SHaRe, place the added ties in the candidate channel, and SHaRe keeps its own similarity-based edge weights in both channels. Two controls test whether a gain requires selecting ties: \textbf{Uniform} runs LightGCN-S with all social entries multiplied by the mean mask weight learned by GBSR, and \textbf{GBSR-Detach} trains GBSR with its mask excluded from gradient updates, as in the public implementation,\footnote{In \texttt{torch\_version/GBSR.py} of \url{https://github.com/yimutianyang/KDD24-GBSR}, the mask weights are detached before the masked adjacency is built.} which keeps the HSIC regularizer but removes learned tie selection.

\textbf{Structure metrics.} On synthetic data we compare the estimates $\hat{c}$, $\hat{\pi}$, $\hat{\rho}^{\text{obs}}$, and $\hat{\rho}^{\text{un}}$ with their true values on $\Omega$ for the training graph (so the true label frequency is $0.9c$), for the Elkan-Noto and the anchored estimator. We measure false-tie detection as the AUROC of the mask $\alpha$ for separating false from genuine observed ties, for both PU-GIB and GBSR. We measure missing-tie recovery in two steps: the fraction of held-out ties $\mathcal{P}_{\text{rec}}$ (and, on synthetic data, of missing edges of $Z$) that enter the candidate pool, against the random base rate and against pools ranked by Resource Allocation, and the AUROC of $\alpha$ for ranking these ties above the other candidates; on real data the positives are the held-out ties of $\mathcal{P}_{\text{rec}}$ in the pool, and pairs of $\mathcal{P}_{\text{est}}$ are excluded from the negatives. A threshold such as $\alpha_{ij} > 0.5$ would be uninformative, because the prior keeps candidate weights near $\hat{\rho}^{\text{un}}$. This also evaluates PU-GIB as a friend recommender, a task GBSR cannot perform. On synthetic data we further report three oracle ceilings: LightGCN-S on the observed graph restricted to genuine ties (denoising), extended by the genuine ties of the candidate pool (addition), and on $Z$ itself.

\textbf{Subgroups.} We report NDCG@20 separately for users with training social degree $0$, $1$--$2$, and at least $3$, and for the bottom quartile of users by number of training interactions. Users without ties receive neither pruning nor candidates, in GBSR and PU-GIB alike, so they serve as a control group, while added ties are expected to matter most for sparse users with few ties.

\textbf{Methods.} \textbf{LightGCN}~\cite{10.1145/3397271.3401063} (no social graph), \textbf{LightGCN-S} (the backbone on the unrefined graph), \textbf{GBSR}~\cite{10.1145/3637528.3671807} (base paper, with a trainable mask as described in the paper), \textbf{SHaRe}~\cite{10.1145/3589334.3645460} (its rewiring rule and contrastive calibration on the same backbone), \textbf{LightGCN-PU} (candidates in the candidate channel with $\tilde{Z} \equiv 1$), and the controls above. Ablations of PU-GIB: \textbf{PU-GIB-Prune} uses observed ties only and isolates the effect of the prior on pruning; a $2 \times 2$ design crosses candidate ranking ($s_{ij}$ or $r_{ij}$) with the prior (estimated or uninformed), giving \textbf{PU-GIB}, \textbf{PU-GIB-Fix}, \textbf{PU-GIB-RA-Est}, and \textbf{PU-GIB-RA}; and \textbf{PU-GIB-Grid} tunes $(\rho^{\text{obs}}, \rho^{\text{un}})$ directly on validation. If PU-GIB-Grid matches PU-GIB, calibration mainly replaces two hyperparameters by one interpretable noise rate. Ablations reuse the hyperparameters selected for PU-GIB.

\textbf{Noise robustness (RQ4).} On the real datasets we inject random false edges amounting to $40\%$ of the training social edges, sampled uniformly from unlinked pairs, recompute $\Omega$ from the corrupted graph, and report the drop in NDCG@20 and the AUROC of the mask for separating injected from genuine ties; no method receives information about the corruption. On synthetic data we use the noise levels listed above.

\textbf{Hyperparameters.} We use $d = 64$, batch size 2048, $L = 3$, Adam with learning rate $10^{-3}$, an $L_2$ weight selected from $\{10^{-5}, 10^{-4}, 10^{-3}\}$ on LightGCN, at most 1500 epochs, and early stopping with a patience of 10 evaluations. The HSIC weight and kernel width follow the GBSR paper on Douban-Book ($\lambda = 40$, $\sigma = 2.5$) and are tuned on validation for the other datasets. The social weights are selected from $\omega \in \{0.5, 1, 1.5, 2, 3\}$ on LightGCN-S and $\omega_{\text{c}} \in \{0.05, 0.1, 0.25\}$ on LightGCN-PU. PU-GIB uses a two-layer mask with 64 hidden units, $K = 10$, $\tau$ annealed from 1.0 to 0.2, $\beta \in \{0, 10^{-3}, 10^{-2}\}$, and $\kappa \in \{0.02, 0.05, 0.1\}$; PU-GIB-Grid searches $\rho^{\text{obs}} \in \{0.7, 0.8, 0.9\}$ and $\rho^{\text{un}} \in \{0.01, 0.05, 0.1\}$.

\subsection{Results}

\section{Conclusion}
% Limitations to discuss here: SCAR within Omega is a population-level
% approximation (SAR extension: label frequency as a function of r_ij);
% the two-hop pool cannot reach missing ties between users without a common
% friend; prior and candidate pool are frozen before training; building
% Omega is quadratic in node degree.

\bibliographystyle{ACM-Reference-Format}
\bibliography{bakhtrob}

\end{document}
